\begin{mistake}{2026-04-15}{01}

\begin{problemcontent}
已知某三阶实对称矩阵的二重特征值对应的特征平面方程为 \(x_1 - 2x_2 + x_3 = 0\)。求属于该特征值的两个正交特征向量（要求求解时第一个特征向量不含有 0 分量）。
\end{problemcontent}

\begin{solutioncontent}
在特征平面方程 \(x_1 - 2x_2 + x_3 = 0\) 中，令 \(x_1=1, x_2=1\)，代入得 \(x_3=1\)。
取第一个特征向量为 \(\xi_1 = [1, 1, 1]^{\mathrm{T}}\)。

设第二个特征向量为 \(\xi_2 = [x_1, x_2, x_3]^{\mathrm{T}}\)。
\(\xi_2\) 需同时满足在特征平面内且与 \(\xi_1\) 正交（即内积 \((\xi_1, \xi_2) = x_1 + x_2 + x_3 = 0\)）。
联立齐次线性方程组：
\[
\begin{cases}
x_1 - 2x_2 + x_3 = 0 \\
x_1 + x_2 + x_3 = 0
\end{cases}
\]
两式相减（下式减上式）得 \(3x_2 = 0 \implies x_2 = 0\)。
代入第二式得 \(x_1 + x_3 = 0 \implies x_3 = -x_1\)。

取自由变量 \(x_1 = 1\)，则 \(x_3 = -1\)。
得到第二个正交特征向量 \(\xi_2 = [1, 0, -1]^{\mathrm{T}}\)。

故所求的两个正交特征向量为 \(\xi_1 = [1, 1, 1]^{\mathrm{T}}\) 与 \(\xi_2 = [1, 0, -1]^{\mathrm{T}}\)。
\end{solutioncontent}

\mistakeknowledge{
\begin{itemize}
 \item \textbf{核心原理}：正交的充要条件为内积为零，即 \((\xi_1, \xi_2)=0\)；求正交特征向量等价于求特征平面与内积正交平面的交线方程解。
 \item \textbf{易错点}：若 \(\xi_1\) 没有 0 分量，无法使用“对应位置互换并变号”的口算捷径凑出 \(\xi_2\)，必须联立并求解方程组，增加计算量和出错率。
 \item \textbf{关联知识}：施密特正交化公式 \(\eta_2 = \alpha_2 - \frac{(\alpha_2, \eta_1)}{(\eta_1, \eta_1)}\eta_1\) 是通法，而联立正交方程法在三阶矩阵二重根的情境下效率更高。
\end{itemize}}

\end{mistake}
